% Test of the headings when keepmaxgap is off (didactic), so that a heading may
% always move: the cases that the limit would otherwise decide.
%
% Build: see the Makefile.  The document checks itself, as test-keep-heading.
%   BA  \section, \subsection, a sentence and a longtable that doesn't fit: both
%       headings move with the table, in the next run, and the outer one isn't
%       left at the foot of the page.
%   BB  a heading, a paragraph of five lines running into a longtable without a
%       blank line: the lines are there, so the heading stays.
%   BC  a heading, two short paragraphs and a longtable: not enough lines before
%       the table, so all move.
%   BD  a heading of two lines, which the estimate takes for less, and a long
%       paragraph: three lines of it with the heading.
%   BB2 the same with three lines, which are enough too.
%   BF  a \chapter, a \section, a sentence and a longtable with a tall row: the
%       chapter isn't left behind by a break before the section.
%   BG  a run-in heading, a line and a tall item that doesn't fit: the heading
%       stays with one line, with a warning (it needs the heading's page).
%   BH  a heading of two lines, a line and a paragraph: the next lines are chained
%       to it, so all move.
%   AQ  two files of the same name in different directories, and one file read
%       twice, the second of each with less room than the first, so that a break
%       recorded for the first would take it too.
%   BE  such a heading, a line and a list, where \LaTeX\ offers a place to break:
%       it takes that within the list, and a warning at the end says which heading
%       has too few lines.
\documentclass[a4paper]{memoir}
\usepackage[british]{babel}
\usepackage{amsmath,amssymb,amsthm}
\usepackage[keepmaxgap=0]{didactic}
\usepackage[noamsthm,notheorems]{beamerarticle}
\usepackage{refcount}
\usepackage{longtable}
\raggedbottom

\makeatletter
\newcommand\roomfor[1]{%
  \clearpage
  \noindent\rule{\linewidth}{0.4pt}\par
  \vspace*{\dimexpr\textheight-\topskip-#1\baselineskip-\baselineskip\relax}%
  \par}
\newcommand\testtable[1]{%
  \begin{longtable}{ll}
    \caption{A caption of several lines, several lines, several lines,
      several lines, several lines, several lines, several lines, several
      lines, several lines.}\label{#1}\\
    a & b\\ a & b\\ a & b\\ a & b\\
  \end{longtable}}
\newcommand\testcheck[4]{%
  \ifnum\getpagerefnumber{#2}#3\getpagerefnumber{#4}\relax
    \def\@tempa{PASS}%
  \else
    \def\@tempa{FAIL}%
  \fi
  \typeout{didactic heading test: \@tempa\space #1 (#2 on \getpagerefnumber{#2}, #4 on \getpagerefnumber{#4})}%
  \item #1: \@tempa}
\makeatother

\begin{document}

\roomfor{8}
\section*{BA section}\label{ba-a}
\subsection*{BA subsection}\label{ba-b}
Table lists the rows.
\testtable{ba-t}
Text after BA.

\roomfor{10}
\section*{BB}\label{bb-a}
Line one of a five line paragraph and so on and so on and so on and on and on.
Line two of a five line paragraph and so on and so on and so on and on and on.
Line three of a five line paragraph and so on and so on and so on and on and on.
Line four of a five line paragraph and so on and so on and so on and on and on.
Line five of a five line paragraph and so on and so on and so on and on and on.\label{bb-5}
\testtable{bb-t}
Text after BB.

\roomfor{9}
\section*{BC}\label{bc-a}
One.

Two.\label{bc-2}
\testtable{bc-t}
Text after BC.

\roomfor{7}
\section*{BD, a heading long enough to run over two lines of its large font,
  indeed it is}\label{bd-a}
First line\label{bd-1}\\
second line\label{bd-2}\\
third line\label{bd-3}\\
fourth line.

\roomfor{6.8}
\section*{BE, a heading long enough to run over two lines of its large font,
  indeed it is}\label{be-a}
One line.
\begin{itemize}
\item An item, which a page may be ended before.\label{be-i}
\item Another item.\label{be-j}
\end{itemize}

\roomfor{8}
\section*{BB2}\label{bbb-a}
Line one of a three line paragraph and so on and so on and on and on and on.
Line two of a three line paragraph and so on and so on and on and on and on.
Line three of a three line paragraph and so on and so on and on and on.\label{bbb-3}
\testtable{bbb-t}
Text after BB2.

\clearpage
\chapter*{BF chapter}\label{bf-a}
\section*{BF section}\label{bf-b}
Table lists the rows.
\begin{longtable}{ll}
\rule{0pt}{.85\textheight}a & b\label{bf-t}\\
\end{longtable}
Text after BF.

\roomfor{4}\label{bg-r}
\paragraph*{BG}\label{bg-a}
One line.
\begin{itemize}
\item\rule{0pt}{5\baselineskip}A tall item.\label{bg-i}
\end{itemize}

\roomfor{6.2}
\section*{BH, a heading long enough to run over two lines of its large font,
  indeed it is}\label{bh-a}
One line.

First line\\
second line\label{bh-2}\\
third line\\
fourth line.

\newcommand\rep{q1}
\newcommand\reproom{5}
\roomfor{\reproom}\label{\rep-r}
\paragraph*{A file read twice}\label{\rep-a}
Table lists the rows.
\ifdim\reproom pt>4.7pt \testtable{\rep-t}\fi
Text after it.

\renewcommand\rep{q2}
\renewcommand\reproom{4.5}
\roomfor{\reproom}\label{\rep-r}
\paragraph*{A file read twice}\label{\rep-a}
Table lists the rows.
\ifdim\reproom pt>4.7pt \testtable{\rep-t}\fi
Text after it.


\roomfor{5}
\paragraph*{First deck}\label{p-a}
Table lists the rows.
\testtable{p-t}
Text after it.

\roomfor{4.5}\label{p-r}
\paragraph*{Second deck}\label{p-c}
Text of the second deck, on the same line of a file of the same name, with
room, so that nothing should move; it is here for the sake of its lines,
which must fit, as they do on this page.


\clearpage
\section*{Checks}
\begin{itemize}
\testcheck{BA outer heading with the inner}{ba-a}{=}{ba-b}
\testcheck{BA headings with the table}{ba-b}{=}{ba-t}
\testcheck{BB heading stays with its lines}{bb-a}{=}{bb-5}
\testcheck{BB table on the next page}{bb-5}{<}{bb-t}
\testcheck{BB2 heading stays with three lines}{bbb-a}{=}{bbb-3}
\testcheck{BB2 table on the next page}{bbb-3}{<}{bbb-t}
\testcheck{BF chapter with the section}{bf-a}{=}{bf-b}
\testcheck{BG run-in heading stays}{bg-r}{=}{bg-a}
\testcheck{BH heading with the second paragraph}{bh-a}{=}{bh-2}
\testcheck{AQ first file: heading with the table}{p-a}{=}{p-t}
\testcheck{AQ second file: its heading isn't moved}{p-r}{=}{p-c}
\testcheck{AQ read twice: first with its table}{q1-a}{=}{q1-t}
\testcheck{AQ read twice: second isn't moved}{q2-r}{=}{q2-a}
\testcheck{BC heading with the table}{bc-a}{=}{bc-t}
\testcheck{BD heading with its third line}{bd-a}{=}{bd-3}
\testcheck{BE the page breaks in the list}{be-i}{<}{be-j}
\end{itemize}
\typeout{didactic heading test: all checks ran}
\end{document}
